% Auto-generated by experiments/build_qualitative_example.py (--top mode)
\begin{figure*}[!htb]\centering
\medskip\noindent\fbox{\begin{minipage}{0.96\linewidth}\small \textbf{Case 1.} \textit{Dataset:} squad \quad \textit{Model:} mistral \\
\textbf{Question.}~What city did Super Bowl 50 take place in? \\
\textbf{Ground truth.}~\textit{Santa Clara}

\textbf{Baseline} ($f=0.942$): The Super Bowl 50 took place in Santa Clara, which is part of the San Francisco Bay Area. \\ 
\textbf{HCPC-v1} ($f=0.478$): San Francisco Bay Area \\ 
\textbf{HCPC-v2} ($f=0.942$, recovered): The Super Bowl 50 took place in Santa Clara, which is part of the San Francisco Bay Area.\end{minipage}}


\medskip\noindent\fbox{\begin{minipage}{0.96\linewidth}\small \textbf{Case 2.} \textit{Dataset:} squad \quad \textit{Model:} mistral \\
\textbf{Question.}~What city did Super Bowl 50 take place in? \\
\textbf{Ground truth.}~\textit{Santa Clara}

\textbf{Baseline} ($f=0.942$): The Super Bowl 50 took place in Santa Clara, which is part of the San Francisco Bay Area. \\ 
\textbf{HCPC-v1} ($f=0.478$): San Francisco Bay Area \\ 
\textbf{HCPC-v2} ($f=0.942$, recovered): The Super Bowl 50 took place in Santa Clara, which is part of the San Francisco Bay Area.\end{minipage}}


\medskip\noindent\fbox{\begin{minipage}{0.96\linewidth}\small \textbf{Case 3.} \textit{Dataset:} squad \quad \textit{Model:} mistral \\
\textbf{Question.}~What year did the Denver Broncos secure a Super Bowl title for the third time? \\
\textbf{Ground truth.}~\textit{2015}

\textbf{Baseline} ($f=0.795$): The provided context does not specify the exact year when the Denver Broncos secured a Super Bowl title for the third time. However, it is implied that this event happened in Super Bowl 50, which was held in 2016. \\ 
\textbf{HCPC-v1} ($f=0.485$): I cannot find this information in the provided context. The context only mentions that the Denver Broncos won the AFC Championship Game and secured home-field advantage throughout the AFC playoffs, but it does not speci\textbackslash\{\}dots \\ 
\textbf{HCPC-v2} ($f=0.795$, recovered): The provided context does not specify the exact year when the Denver Broncos secured a Super Bowl title for the third time. However, it is implied that this event happened in Super Bowl 50, which was held in 2016.\end{minipage}}


\medskip\noindent\fbox{\begin{minipage}{0.96\linewidth}\small \textbf{Case 4.} \textit{Dataset:} squad \quad \textit{Model:} llama-3.3-70b \\
\textbf{Question.}~If Roman numerals were used, what would Super Bowl 50 have been called? \\
\textbf{Ground truth.}~\textit{Super Bowl L}

\textbf{Baseline} ($f=0.502$): If Roman numerals were used, Super Bowl 50 would have been called Super Bowl L. \\ 
\textbf{HCPC-v1} ($f=0.491$): Super Bowl L. \\ 
\textbf{HCPC-v2} ($f=0.502$, recovered): If Roman numerals were used, Super Bowl 50 would have been called Super Bowl L.\end{minipage}}

\caption{\textbf{Real-world case studies of the refinement paradox.} Five queries where the same question is answered correctly with raw retrieval, incorrectly after per-passage refinement, and correctly again under coherence-gated refinement. Faithfulness $f$ is the NLI entailment score (higher is better; $\geq 0.5$ is the conventional faithful threshold). All cases satisfy $\textrm{baseline} = \textrm{faithful}$, $\textrm{HCPC-v1} = \textrm{hallucinated}$, $\textrm{HCPC-v2} = \textrm{faithful}$.}
\label{fig:qualitative_cases}
\end{figure*}
